\section{Guarded predicating: discharging the proof obligation}

\Cref{thm:predicating-proper} states that predicating is defined on an
articulation cell only with a proof that the cell is in the potential corner,
but the carrier it is stated against, \texttt{Articulating.PotentialCell}, is
an unconstrained wrapper: its ``potential''-ness is asserted by the type's
name, not demanded by a hypothesis. This section discharges the promise
literally, against a general \texttt{Articulating.Cell} with an explicit
incidence proof as an argument, and records the resulting carrier as
equivalent to the existing typed one.

\begin{proposition}[The typed potential cell]\label{thm:potential-cell}\lean{Obsmetrologia.Articulating.PotentialCell}\leanok\uses{thm:articulation-cell}
The typed potential cell is an utterance alone, read as told-true (token
witnessed, claim not) by construction: \texttt{toCell} forgets it to the
ambient \texttt{Articulating.Cell} via \texttt{articulate}, giving token slot
\texttt{[utterance.token]} and empty claim slot, and \texttt{toCell\_corner}
confirms the forgotten cell always sits in the \texttt{potential} corner. The
incidence fact is true of every value of this type, but nothing in the type
itself forces a caller to supply it.
\end{proposition}

\begin{proposition}[Predicating, guarded]\label{thm:predicating-guarded}\lean{Obsmetrologia.Articulating.Cell.predicatingOf}\leanok\uses{thm:predicating-proper, thm:corner}
\texttt{predicatingOf} takes a general articulation cell \(c\) together with
an explicit proof \(c.\mathrm{corner} = \mathtt{potential}\) as its second
argument, and returns the seed of \(c\)'s utterance. Computationally this is
the same construction \texttt{PredicationCell.predicating} always was; the
proof argument is not consumed in the body. What changes is the type: the
incidence condition the prose promised is now a hypothesis the caller must
supply, not a name the caller must trust.
\end{proposition}

\begin{theorem}[The guarded and typed carriers agree]\label{thm:predicating-guarded-agree}\lean{Obsmetrologia.PredicationCell.predicatingOf_toCell}\leanok\uses{thm:predicating-guarded, thm:potential-cell, thm:predicating-proper}
Guarded predicating on the cell \(pc.\mathrm{toCell}\) underlying a typed
potential cell \(pc\), fed the proof \texttt{toCell\_corner} already supplies,
equals \texttt{predicating}~\(pc\). Both sides reduce to
\texttt{seed}~\(pc.\mathrm{utterance}\); the equality is definitional.
\end{theorem}

\begin{proposition}[The potential-corner subtype]\label{thm:potential-subtype}\lean{Obsmetrologia.Articulating.PotentialCell.Potential}\leanok\uses{thm:corner, thm:articulation-cell-corner}
The second carrier of ``a cell known to sit in the potential corner'': the
subtype of articulation cells \(c\) paired with a proof
\(c.\mathrm{corner} = \mathtt{potential}\). This is the general-purpose,
Mathlib-shaped presentation of the same incidence condition
\texttt{PotentialCell} carries by convention rather than by type.
\end{proposition}

\begin{theorem}[The two carriers are a retract, not a bijection]\label{thm:potential-roundtrip}\lean{Obsmetrologia.Articulating.PotentialCell.ofSubtype_toSubtype, Obsmetrologia.Articulating.PotentialCell.toSubtype_ofSubtype}\leanok\uses{thm:potential-cell, thm:potential-subtype}
\texttt{toSubtype} carries a typed potential cell to the potential-corner
subtype, and \texttt{ofSubtype} forgets a subtype element's incidence proof
back to a typed potential cell, keeping only its utterance. Under the slot
cell, \texttt{PotentialCell} is a \emph{retract} of the subtype, not an
equivalence with it: forward-then-back,
\texttt{ofSubtype\,(toSubtype\,pc) = pc}, is unconditional and definitional
(\texttt{ofSubtype\_toSubtype}). Back-then-forward,
\texttt{toSubtype\,(ofSubtype\,c) = c}, needs the extra hypothesis
\(c.\mathrm{val}.\mathrm{tokenSlot} = [c.\mathrm{val}.\mathrm{utterance}.\mathrm{token}]\)
(\texttt{toSubtype\_ofSubtype}): the incidence proof
\(c.\mathrm{val}.\mathrm{corner} = \mathtt{potential}\) alone only pins the
claim slot empty and the token slot non-empty, not to the singleton
\([c.\mathrm{val}.\mathrm{utterance}.\mathrm{token}]\) --- a potential cell may
carry further token witnesses beyond the one \texttt{articulate} would have
put there, and \texttt{ofSubtype} (which keeps only the utterance) silently
forgets them. Without the hypothesis the round trip is false in general; with
it, the two carriers of ``told-true'' coincide on the canonical, measured
shape.
\end{theorem}

\begin{proposition}[What always survives the round trip]\label{thm:potential-roundtrip-survives}\lean{Obsmetrologia.Articulating.PotentialCell.toSubtype_ofSubtype_utterance, Obsmetrologia.Articulating.PotentialCell.toSubtype_ofSubtype_corner}\leanok\uses{thm:potential-roundtrip}
Even without the singleton hypothesis, \texttt{toSubtype\,(ofSubtype\,c)}
always keeps \(c\)'s own utterance and its \texttt{potential} corner: the
back-then-forward composite is unconditionally a potential cell over the same
utterance as \(c\), and only the extra token witnesses beyond the head's own
token can be lost. This is the honest unconditional remainder of the full
round trip \cref{thm:potential-roundtrip} loses in general.
\end{proposition}

\begin{proposition}[A cell at the un-articulated position cannot predicate]\label{thm:absurd-cannot-predicate}\lean{Obsmetrologia.Articulating.Cell.absurd_cannot_predicate}\leanok\uses{thm:predicating-guarded, thm:articulation-cell-initial}
A cell in the \texttt{absurd} corner --- the un-articulated position --- can
never supply \texttt{predicatingOf}'s hypothesis: \texttt{absurd} and
\texttt{potential} are distinct constructors of \texttt{Corner}. This is
deliberately shallow --- a naming of constructor-distinctness, not a deep
result --- but it is the honest Lean shape of the qualitative claim that
predicating cannot proceed from the un-articulated position: at least one
articulation, witnessing the token, is required first. Since articulating
itself never produces the absurd corner (\cref{thm:cell-corner-ne-absurd}),
this hypothesis is idle against anything articulating actually generates; it
guards only against a cell built by fiat.
\end{proposition}

This section leaves \texttt{PredicationCell.predicating} itself untouched;
it does not retrofit the guard onto the carrier the rest of the predicating
chapter is built against, nor does it claim the retract of
\cref{thm:potential-roundtrip} is the unique such relation between the two
carriers of ``potential'' --- only that this one exists, with its hypothesis
made explicit rather than smoothed over.
